% GENERATED FILE. Edit canonical lessons, then run npm run generate:books.
% canonical-source-hash: fnv1a64:3b89ee6dc34a289c
% canonical-lessons: ZH-C143-miao, ZH-C143-fenzhong, ZH-C143-cha, ZH-C143-zhunshi, ZH-C143-bantian

\chapter{Clock Words}
\label{ch:zh-clock-words}
\hlchaptermodality{143}

\begin{chapteropening}
\textbf{By the end of this chapter:} \emph{I can say seconds and minutes, the minutes before the hour, and when things begin and end.}

Complete the last lesson of chapter 143: I can say seconds and minutes, the minutes before the hour, and when things begin and end.
\end{chapteropening}

\section[miǎo]{miǎo — second (of time)}

\label{lesson:ZH-C143-miao}



\begin{quote}
Before the new one: say the Mandarin for to sort out, then the Mandarin for to water the plants.
\end{quote}

\subsection*{You'll want to know: miǎo}
\emph{miǎo} — \textquotedblleft{}second (of time)\textquotedblright{}.

\begin{quote}
\emph{qǐng děng wǔ miǎo} — Please wait five seconds.
\end{quote}

Its characters come later; for now it is a word for the ear.

\begin{cousinweb}
\emph{miǎo} is inherited from Middle Chinese \emph{mjiewX} (Guangyun, in Baxter's transcription), glossed 'the beard of grain'. The tiny unit of time is not given there; the step from something very fine is the likely one, but it is not settled.
\end{cousinweb}

\begin{grammarlens}[title={What you've built}]
One more word for times of day.
\end{grammarlens}

\subsection*{Your turn}
\begin{itemize}
\raggedright
  \item \emph{Say it:} \emph{miǎo}
  \item \emph{Say it:} \emph{miǎo}, once more
  \item \emph{Say it:} \emph{qǐng děng wǔ miǎo} — Please wait five seconds.
  \item \emph{Recall:} say the Mandarin for to sort out, then the Mandarin for to water the plants, then say \emph{miǎo} again
\end{itemize}

\begin{tcolorbox}[breakable,colback=teal!4,colframe=teal!35!black,title={Before you move on}]
What does miǎo mean? (\textquotedblleft{}Second (of time)\textquotedblright{}.) Say it once more.
\end{tcolorbox}
\section[fēnzhōng]{fēnzhōng — minute (a length of time)}

\label{lesson:ZH-C143-fenzhong}



\begin{quote}
Before the new one: say the Mandarin for to water the plants, then the Mandarin for second (of time).
\end{quote}

\subsection*{You'll want to know: fēnzhōng}
\emph{fēnzhōng} — \textquotedblleft{}minute (a length of time)\textquotedblright{}.

\begin{quote}
\emph{shí fēnzhōng yǐhòu wǒmen zǒu} — We're leaving in ten minutes.
\end{quote}

Its characters come later; for now it is a word for the ear.

\begin{cousinweb}
\emph{fēnzhōng} is built from \emph{fēn}, the minute of the clock, and \emph{zhōng} 'bell', both inherited from Middle Chinese, \emph{pjun} and \emph{tsyowng} (Guangyun, in Baxter's transcription); the dictionary glosses the second as a musical instrument.
\end{cousinweb}

\begin{grammarlens}[title={What you've built}]
One more word for times of day.
\end{grammarlens}

\subsection*{Your turn}
\begin{itemize}
\raggedright
  \item \emph{Say it:} \emph{fēnzhōng}
  \item \emph{Say it:} \emph{fēnzhōng}, once more
  \item \emph{Say it:} \emph{shí fēnzhōng yǐhòu wǒmen zǒu} — We're leaving in ten minutes.
  \item \emph{Recall:} say the Mandarin for to water the plants, then the Mandarin for second (of time), then say \emph{fēnzhōng} again
\end{itemize}

\begin{tcolorbox}[breakable,colback=teal!4,colframe=teal!35!black,title={Before you move on}]
What does fēnzhōng mean? (\textquotedblleft{}Minute (a length of time)\textquotedblright{}.) Say it once more.
\end{tcolorbox}
\section[chà]{chà — to be short of; to (before the hour)}

\label{lesson:ZH-C143-cha}



\begin{quote}
Before the new one: say the Mandarin for second (of time), then the Mandarin for minute (a length of time).
\end{quote}

\subsection*{You'll want to know: chà}
\emph{chà} — \textquotedblleft{}to be short of; to (before the hour)\textquotedblright{}.

\begin{quote}
\emph{xiànzài chà wǔ fēn bā diǎn} — It's five to eight now.
\end{quote}

Its characters come later; for now it is a word for the ear.

\begin{cousinweb}
\emph{chà} is inherited from Middle Chinese, where the Guangyun lists several readings; \emph{tsrhae} (Guangyun, in Baxter's transcription) is glossed 'to choose' and 'amiss'. None of them gives the fourth tone of \emph{chà} by the regular sound change, so how this reading came is not settled.
\end{cousinweb}

\begin{grammarlens}[title={What you've built}]
One more word for times of day.
\end{grammarlens}

\subsection*{Your turn}
\begin{itemize}
\raggedright
  \item \emph{Say it:} \emph{chà}
  \item \emph{Say it:} \emph{chà}, once more
  \item \emph{Say it:} \emph{xiànzài chà wǔ fēn bā diǎn} — It's five to eight now.
  \item \emph{Recall:} say the Mandarin for second (of time), then the Mandarin for minute (a length of time), then say \emph{chà} again
\end{itemize}

\begin{tcolorbox}[breakable,colback=teal!4,colframe=teal!35!black,title={Before you move on}]
What does chà mean? (\textquotedblleft{}To be short of; to (before the hour)\textquotedblright{}.) Say it once more.
\end{tcolorbox}
\section[zhǔnshí]{zhǔnshí — on time; punctual}

\label{lesson:ZH-C143-zhunshi}



\begin{quote}
Before the new one: say the Mandarin for minute (a length of time), then the Mandarin for to be short of.
\end{quote}

\subsection*{You'll want to know: zhǔnshí}
\emph{zhǔnshí} — \textquotedblleft{}on time; punctual\textquotedblright{}.

\begin{quote}
\emph{huǒchē hěn zhǔnshí} — The trains run on time.
\end{quote}

Its characters come later; for now it is a word for the ear.

\begin{cousinweb}
\emph{zhǔnshí} is built from \emph{zhǔn} 'level, a standard', the \emph{zhǔn} of \emph{zhǔnbèi}, and \zh{时} 'time', both inherited from Middle Chinese, \emph{tsywinX} and \emph{dzyi} (Guangyun, in Baxter's transcription): true to the time.
\end{cousinweb}

\begin{grammarlens}[title={What you've built}]
One more word for times of day.
\end{grammarlens}

\subsection*{Your turn}
\begin{itemize}
\raggedright
  \item \emph{Say it:} \emph{zhǔnshí}
  \item \emph{Say it:} \emph{zhǔnshí}, once more
  \item \emph{Say it:} \emph{huǒchē hěn zhǔnshí} — The trains run on time.
  \item \emph{Recall:} say the Mandarin for minute (a length of time), then the Mandarin for to be short of, then say \emph{zhǔnshí} again
\end{itemize}

\begin{tcolorbox}[breakable,colback=teal!4,colframe=teal!35!black,title={Before you move on}]
What does zhǔnshí mean? (\textquotedblleft{}On time; punctual\textquotedblright{}.) Say it once more.
\end{tcolorbox}
\section[\texorpdfstring{\zh{半天} (bàntiān)}{半天 (bàntiān)}]{\zh{半天} (bàntiān) — a long while; half a day}

\label{lesson:ZH-C143-bantian}



\begin{quote}
Before the new one: say the Mandarin for to be short of, then the Mandarin for punctual.
\end{quote}

\subsection*{You'll want to know: \zh{半天}}
\textbf{\zh{半天}} — \emph{bàntiān} — \textquotedblleft{}a long while; half a day\textquotedblright{}.

\begin{quote}
\emph{wǒ děng le bàntiān} — I waited for ages.
\end{quote}

\begin{cousinweb}
\textbf{\zh{半天}} (\emph{bàntiān}) is built from \zh{半} 'half' and \zh{天} 'day', both inherited from Middle Chinese, \emph{panH} and \emph{then} (Guangyun, in Baxter's transcription): half a day, and in everyday talk, often with a sigh, a long while.
\end{cousinweb}

\begin{grammarlens}[title={What you've built}]
One more word for times of day.
\end{grammarlens}

\subsection*{Your turn}
\begin{itemize}
\raggedright
  \item \emph{Say it:} \emph{bàntiān}
  \item \emph{Say it:} \emph{bàntiān}, once more
  \item \emph{Say it:} \emph{wǒ děng le bàntiān} — I waited for ages.
  \item \emph{Recall:} say the Mandarin for to be short of, then the Mandarin for punctual, then say \emph{bàntiān} again
\end{itemize}

\begin{tcolorbox}[breakable,colback=teal!4,colframe=teal!35!black,title={Before you move on}]
What does \zh{半天} mean? (\textquotedblleft{}A long while; half a day\textquotedblright{}.) Say it once more.
\end{tcolorbox}
